% !TeX root = ../初等数学笔记.tex
% 由主文档通过 \input 载入；请编译主文档。
\setcounter{chapter}{7} % 保留原书章号。
\chapter{幂、指数与对数}
\label{chap:em-exponential}
线性模型描述每一步增加固定数量，指数模型描述每一步乘以固定比例。对数则反过来回答“需要多少次比例变化才能达到目标”。二者互逆，是处理尺度和增长的基本工具。

\section{有理指数与实指数的扩展}
对 $a>0$、整数 $p$ 和正整数 $q$，定义 $a^{p/q}=\sqrt[q]{a^p}$。若 $p/q=r/s$，把两个正数的幂提升到 $qs$ 次，得到同样的 $a^{ps}=a^{rq}$，由正数幂的单射性知定义不依赖分数表示。利用同样方法可证明有理指数的加法与乘法法则。

对 $a>1$ 和实数 $x$，可以利用完备性定义
\[
 a^x=\sup\{a^r:r\in\mathbb Q,\ r\le x\}.
\]
集合非空且有上界：选择整数 $m\le x<M$，则 $a^m$ 是其中元素，$a^M$ 是上界。对 $0<a<1$ 定义 $a^x=(1/a)^{-x}$，对 $a=1$ 定义 $1^x=1$。这一构造与有理指数相容，并给出连续、保持指数律的唯一延拓；完整证明还需实数逼近理论，本章采用该延拓性质，不以有限位小数试算代替证明。底数为正是这些实指数法则的共同前提。

\section{幂函数与定义域}
\emph{幂函数}（power function）把指数固定，如 $x^2,x^{-1},x^{1/2}$。整数正幂在实轴上都有定义，负整数幂排除零；偶次根限制输入非负，奇次根可接受负数。一般实指数的 $x^\alpha$ 默认在 $x>0$ 上讨论，再按具体指数判断能否扩展到零或负数。

对 $x>0$，$\alpha>0$ 时 $x^\alpha$ 严格递增，$\alpha<0$ 时严格递减，$\alpha=0$ 时恒为 $1$。整数幂可以用因式分解直接证明，例如 $0<u<v$ 时 $v^n-u^n=(v-u)(v^{n-1}+\cdots+u^{n-1})>0$；实指数情形依赖上一节的延拓性质。实负底数的分数幂则须单独定义，不能任意交换幂次顺序，例如 $((-1)^2)^{1/2}=1$ 而 $(-1)^{2\cdot(1/2)}=-1$。

\section{指数函数与增长模型}
固定 $a>0,a\ne1$ 时，$f(x)=a^x$ 称为\emph{指数函数}（exponential function）。它的定义域为 $\mathbb R$，值域为 $(0,+\infty)$，经过 $(0,1)$；$a>1$ 时严格递增，$0<a<1$ 时严格递减，并满足
\[
 a^{x+y}=a^xa^y,\qquad a^{x-y}=\frac{a^x}{a^y}.
\]
底数大小决定相同时间间隔的比例变化。例如初值 $V_0>0$、每期增长率 $r>-1$ 恒定时，离散模型为 $V_n=V_0(1+r)^n$。把整数期数延拓为实变量是额外建模选择，并不意味着实际过程在每一瞬间都遵循同一曲线。

\section{对数、换底与数量级}
\begin{definition}
设 $a>0,a\ne1,x>0$。满足 $a^y=x$ 的唯一实数 $y$ 称为 $x$ 以 $a$ 为底的\emph{对数}（logarithm），记为 $\log_a x$。
\end{definition}
指数函数连续、严格单调且覆盖全部正数，保证定义中的存在与唯一。于是 $\log_a1=0,\log_a a=1$，以及
\begin{equation}\label{eq:em-loglaws}
 \log_a(xy)=\log_a x+\log_a y,\quad
 \log_a(x/y)=\log_a x-\log_a y,\quad
 \log_a(x^t)=t\log_a x
\end{equation}
对 $x,y>0,t\in\mathbb R$ 成立。例如令 $u=\log_a x,v=\log_a y$，则 $xy=a^{u+v}$，由唯一性得到第一式，其他两式同理。

若 $b>0,b\ne1$，则 $a^{\log_a x}=x$ 两侧取 $\log_b$ 给出换底公式 $\log_a x=\log_b x/\log_b a$。底数 $10$ 的对数反映数量级；底数为常数 $e$ 的对数记为 $\ln x$，$e$ 的一种构造见第\ref{chap:em-modeling}章。对数不把加法变为加法，例如 $\log_{10}(1+1)>0$，而 $\log_{10}1+\log_{10}1=0$。

\section{指数与对数方程、不等式}
指数函数单射，故 $a^{u(x)}=a^{v(x)}$ 等价于 $u(x)=v(x)$，前提是固定底数合法且两指数有定义。指数不等式的方向由底数决定；对数方程则先限制真数为正。

例如 $\log_2(x-1)+\log_2(x+1)=3$ 的定义域为 $x>1$，由式 \eqref{eq:em-loglaws} 化为 $x^2-1=8$，两候选值 $\pm3$ 中只留下 $3$。又 $\log_{1/2}(x-1)\ge2$ 先要求 $x>1$；由于底数小于 $1$，得到 $0<x-1\le1/4$，故 $1<x\le5/4$。忽略单调方向会得到完全相反的区间。

\section{例题、习题与学习检查}
\begin{example}
一笔自拟资金按每年 $5\%$ 的固定比例增长，至少经过多少个完整年度才翻倍？
\end{example}
\begin{solve}
要求 $1.05^n\ge2$，$n$ 为非负整数，所以 $n\ge\log_{1.05}2$。直接计算可验证 $1.05^{14}<2<1.05^{15}$，故最小完整年数为 $15$。该答案依赖增长率恒定且收益持续计入本金的假设，不是对实际收益的预测。
\end{solve}
\begin{enumerate}
\item 解 $3^{2x-1}=27$。答：$x=2$。
\item 解 $\log_2(x+2)-\log_2x=1$。答：$x=2$，定义域 $x>0$。
\item 解 $(1/3)^{x-1}<9$。答：$x>-1$。
\end{enumerate}

